Studiamo lo spettro di emissione di raggi X di un tubo Röntgen.
Misurando la radiazione con un contatore Geiger-Müller (di cui caratterizziamo il funzionamento) verifichiamo la legge di Lambert-Beer stimando il coefficiente di assorbimento dell'alluminio $ \mu = (\ldots \pm \ldots) $. 
Poi, analizziamo lo spettro di emissione del tubo radiogeno studiandone la diffrazione su un cristallo di \ce{NaCl}: le lunghezze d'onda delle due righe di emissione principali sono \ldots
Date queste stime, siamo anche in grado di valutare la distanza interatomica di cristalli di \ldots come \ldots
